%-------------------------
% Cover Letter in Latex
% Enrique Martínez Martel
%------------------------

\documentclass[letterpaper,11pt]{article}

\usepackage{latexsym}
\usepackage[empty]{fullpage}
\usepackage{marvosym}
\usepackage[usenames,dvipsnames]{color}
\usepackage[hidelinks]{hyperref}
\usepackage{fancyhdr}
\usepackage[english]{babel}
\usepackage{tabularx}
\usepackage{ragged2e}
\usepackage{fontawesome5}
\usepackage{graphicx}
\usepackage[scale=0.90,lf]{FiraMono}
\usepackage{enumitem}

\definecolor{light-grey}{gray}{0.83}
\definecolor{dark-grey}{gray}{0.3}
\definecolor{text-grey}{gray}{.08}

\usepackage{contour}
\usepackage[normalem]{ulem}
\renewcommand{\ULdepth}{1.8pt}
\contourlength{0.8pt}
\newcommand{\myuline}[1]{%
  \uline{\phantom{#1}}%
  \llap{\contour{white}{#1}}%
}

\usepackage{tgheros}
\renewcommand*\familydefault{\sfdefault}
\usepackage[T1]{fontenc}

\pagestyle{fancy}
\fancyhf{}
\fancyfoot{}
\renewcommand{\headrulewidth}{0pt}
\renewcommand{\footrulewidth}{0pt}

\addtolength{\oddsidemargin}{-0.5in}
\addtolength{\evensidemargin}{0in}
\addtolength{\textwidth}{1in}
\addtolength{\topmargin}{-.8in}
\addtolength{\textheight}{1.5in}

\urlstyle{same}
\raggedbottom
\raggedright

\color{text-grey}

\begin{document}

\vspace*{160pt}

Dear Professor Vandergheynst,

\vspace{16pt}

I am applying for the PhD position in Multimodal Generative AI for Materials Design at the LTS2 Lab (Reference 2048). My background is in deep learning and computer vision: an MSc thesis at UPM benchmarking YOLOv10 and Faster R-CNN on real-world data (grade: 10/10), a multimodal LLM classification pipeline built at Banco Santander, and a BSc thesis on autonomous driving perception at UPC. More of my work is at \href{https://enkairito.github.io/enriquemm.github.io/}{\myuline{enkairito.github.io/enriquemm.github.io}}.

\vspace{16pt}

The main reasons I want to do this PhD:

\vspace{8pt}

\begin{itemize}[leftmargin=1.5em, itemsep=10pt]

  \item \textbf{The closed-loop problem.} Most of the ML work I have done has been evaluated offline. The idea of a generative model that feeds into an automated synthesis platform and gets corrected by actual experimental results is a fundamentally different kind of problem — and one I find much more interesting.

  \item \textbf{Where domain knowledge shapes the model.} The project forces you to understand synthesis conditions, characterisation data, and simulation outputs — not as labels, but as constraints on what the model should learn. That intersection between domain expertise and architecture design is where I think the most rigorous research happens, and working alongside materials scientists and physicists to get there is the kind of environment I want to be in.

  \item \textbf{Generative AI applied to something physically real.} Most generative AI work I have seen in industry operates on language or images with no ground truth beyond human preference. Using latent diffusion to navigate the space of possible materials — where the output can be synthesised and tested against reality — is a fundamentally harder and more interesting problem. That is where I want to develop as a researcher.

\end{itemize}

\vspace{16pt}

I do not have a materials science background. I am not pretending otherwise. But I learn new domains fast, and I am drawn to hard problems in new areas.

\vspace{16pt}

I would be happy to discuss further.

\vspace{20pt}

Best regards,

\vspace{16pt}

\textbf{Enrique Martínez Martel}

\vspace{6pt}

\faEnvelope \hspace{2pt} \texttt{enri.martinez25@gmail.com} \hspace{6pt} $|$ \hspace{6pt}
\faLinkedin \hspace{2pt} \href{https://www.linkedin.com/in/enrique-martinez-martel}{\myuline{enrique-martinez-martel}} \hspace{6pt} $|$ \hspace{6pt}
\faGlobe \hspace{2pt} \href{https://enkairito.github.io/enriquemm.github.io/}{\myuline{enkairito.github.io/enriquemm.github.io}}

\end{document}
